% !TeX root = ../教材题库_全部条目_转写工作稿.tex
% Source: XJ-B1-C06-RV-P266-Y0326
% 第6章 / 章末复习 复习题六 / 教材 p.266 / PDF 第273页
% OCR未校对：公式与文字均可能有误；未完成数学TeX转写；禁止直接发给学生
% 原题及相邻原解截图（核对用）：../题目原图/XJ-B1-C06-RV-P266-Y0326-01.webp
% 本文件只在转写完成后作为题干使用；OCR 条目尚不是可靠题干。
\begin{bankitem}{复习题六 第10题}{第6章 / 章末复习 复习题六 / 教材 p.266 / PDF 第273页}{XJ-B1-C06-RV-P266-Y0326}{OCR未校对：公式与文字均可能有误；未完成数学TeX转写；禁止直接发给学生}
\begin{ocrdraft}
[教材p.266 / PDF第273页]
\par
10．为提高科技投入效能，某机构对某“智慧园区”的20家企业的科技成果
\par
转化效能进行评估，获得的评估分数如下：
\par
786372
\par
899156
\par
68
\par
76
\par
85
\par
60
\par
718461
\par
897993
\par
86
\par
78
\par
92
\par
80
\par
（D以上述数据组成总体，
\par
求总体平均数与总体标准差．
\par
（2）设计恰当的随机抽样方法，从总体中抽取一个容量为10的样本，求样本
\par
均值与标准差．
\par
（3）利用上面的随机抽样方法，再抽取容量为10的样本，计算样本均值和标
\par
准差．将求得的结果与（2）中的结果进行比较，它们一样吗？
\par
（4）利用（2）中的随机抽样方法，分别从总体中抽取一个容量为8，12，16，
\par
18的样本，求样本均值与标准差．分析样本容量与样本均值、样本标准差对总体
\par
的估计效果之间有什么关系．
\end{ocrdraft}
\end{bankitem}
